\documentclass[10pt,a4paper]{amsart}
\usepackage[margin=19mm]{geometry}
\usepackage{amsmath,amssymb,mathtools}
\usepackage[scr=rsfs,cal=euler]{mathalfa}
\usepackage{xfrac}
\usepackage[unicode,hidelinks,bookmarks=false]{hyperref}
\newcommand{\ie}{i.e.\ }
\newcommand{\cf}{cf.\ }
\newcommand{\resp}{respectively\ }
\newcommand{\Subsection}{\subsection}
\providecommand{\nobreakdash}{\nobreak\hbox{-}}
\setlength{\emergencystretch}{2em}
\multlinegap0pt
\usepackage{bm}
\usepackage{bbm}
\usepackage{dsfont}
\usepackage{xspace}
\usepackage{cancel}
\usepackage{mathtools}
\usepackage{amsthm}
\makeatletter
\@ifundefined{observation}{\newtheorem{observation}{Observation}}{}
\@ifundefined{theorem}{\newtheorem{theorem}{Theorem}}{}
\makeatother
\title[On a conjecture about the strong odd chromatic number of pla]{On a conjecture about the strong odd chromatic number of planar graphs}
\author{Arun J Manattu, Athira Vinay, Aparna Lakshmanan S}
\date{Source: arXiv:2602.03259v1 (CC BY 4.0)}
\begin{document}
\maketitle

\noindent\emph{Source and licence.} On a conjecture about the strong odd chromatic number of planar graphs. Version arXiv:2602.03259v1 (CC BY 4.0), distributed under the Creative Commons Attribution 4.0 International licence, as recorded in the arXiv licence metadata for this exact version and shown on the abstract page https://arxiv.org/abs/2602.03259v1. Full rights notice: licenses/arxiv-2602.03259-CC-BY-4.0.txt.\par\smallskip

\noindent\textit{Editorial scope.} Counted unit: the Theorem whose source label is \texttt{wheels} in the article (source0.tex lines 121--131, source label \texttt{wheels}), delivered together with the article's own complete original proof of it (source0.tex lines 133--147, one \texttt{proof} environment of 15 lines). No mathematics is written, repaired or re-derived: statement and proof are the article's own text line for line. Required context retained uncounted: H odd (lines 111--113); d>2 (lines 117--119). Every definition, display and label of those lines is retained unchanged. Omitted from this package and not redistributed: 1--110, 114--116, 120--120, 132--132, 148--617. Nothing inside a retained excerpt is omitted or altered: every retained body is delivered exactly as the article's lines are written, so no omission receipt of any kind accompanies this package. Citations: the retained excerpts carry no citation command at all, so no bibliography entry of the article is transcribed and no reference is redistributed. Reference adaptation: none was needed and none was made; every reference inside the delivered unit resolves inside it, so no unresolved reference or citation is produced. Printed slips kept literal: no printed word, symbol, hypothesis or number of the counted statement or of its proof was altered. Layout and class: the article's own class is \texttt{article} with packages including amsmath, amsthm, amscd, amsfonts, amssymb, graphicx, color, accents, hyperref, euscript, latexsym, float, multicol, tabularx; the wrapper is the lane's pinned \texttt{amsart} editorial wrapper at 10pt on A4, which restates the theorem-like environments the retained text needs and adds the article's own macro definitions that the retained text needs (none), with extra wrapper packages (bm, bbm, dsfont, xspace, cancel, mathtools). The delivered numbering is the unit's own. Assets: no retained span contains a figure environment, a graphics-inclusion command, a PDF-page inclusion, a table or a TikZ picture; no image file is copied into this package, the package declares no asset dependency and no image file is redistributed with it. Licence: version arXiv:2602.03259v1 of this article is distributed under the Creative Commons Attribution 4.0 International licence, as recorded in the arXiv licence metadata for this exact version and shown on the abstract page https://arxiv.org/abs/2602.03259v1. A full rights notice is delivered at licenses/arxiv-2602.03259-CC-BY-4.0.txt. Characters: every retained excerpt is pure ASCII, so the delivered engine, XeLaTeX with its default Latin Modern text font, reports no missing character.
\begin{observation} \label{H odd}
    In every strong odd coloring of $C_n \vee \overline{K_m}$, every color appearing in $C_n$ must be an odd color class of $C_n$.
\end{observation}

\begin{observation} \label{d>2}
    A strong odd coloring of $C_n \vee \overline{K_m}$, when restricted to the vertices in $C_n$ gives a 2-distance coloring.
\end{observation}

\begin{theorem} \label{wheels}
Let $W_n = C_n \vee K_1$ be the \textit{wheel graph} on $n+1$ vertices. Then,

    $\chi_{\text{so}}(W_n) = \begin{cases}
  4, & \text{if }~ n \equiv 3(mod~6)\\
  5,  & \text{if }~n \equiv 0, ~2 \text{ or }4(mod~6), ~n>8,~ n\neq15\\
  6,  & \text{if }~ n \equiv 1 \text{ or }5(mod~6), ~n>8\\
  7, & \text{if }~ n = 14\\
  n+1, & \text{if }~ n \leq 8\\  
\end{cases}$
\end{theorem}

\begin{proof}
If $n \leq 8$, by Observation \ref{d>2}, no color can appear more than 2 times in the cycle since among any three vertices of a cycle with at most 8 vertices, there is at least one pair with distance at most 2. But then, by Observation \ref{H odd}, this implies that every color appears exactly once, i.e; all vertices are colored with distinct colors. So, $\chi_{\text{so}}(C_n \vee K_1) = n+1$, for $ n\leq 8$. \par
For $n>8$, as a direct consequence of Observation \ref{d>2}, a color can be used at most $\lfloor \frac{n}{3} \rfloor$ times to color the vertices of $C_n$. Also, by applying Observation \ref{H odd}, a color can be used at most $\lfloor \frac{n}{3} \rfloor$ times, if $\lfloor \frac{n}{3} \rfloor$ is odd and $\lfloor \frac{n}{3} \rfloor - 1$ times, otherwise. So, when $n \equiv 0,1 \text{ or } 2\text{ (mod 6)}$, a color can be used at most $\lfloor \frac{n}{3} \rfloor - 1$ times to color $C_n$ and it can be used at most $\lfloor \frac{n}{3} \rfloor$ times, when $n \equiv 3,4 \text{ or } 5 \text{ (mod 6)}$.\par

When $n \equiv 0 \text{ (mod 6)}$, we need at least 4 colors to color $C_n$. Also, the colors $\{1,2,3,4\}$ can be sequentially used 3 times to color the first 12 vertices and the remaining vertices, if any, can be colored sequentially using the colors $\{1,2,3\}$.\par

When $n \equiv 1 \text{ or }5 \text{ (mod 6)}$, then again, we need at least 4 colors to color $C_n$. But, $n$ being odd and the cardinality of each color class being odd (by Observation \ref{H odd}), the number of colors must also be odd. So, we need at least 5 colors to color $C_n$. Now, for $n \equiv 1 \text{ (mod 6)}$, the colors $\{1,2,3,4\}$ can be sequentially used 3 times to color the first 12 vertices, color 5 to the next vertex and the remaining vertices, if any, can be colored sequentially using the colors $\{1,2,3\}$. For $n \equiv 5 \text{ (mod 6)}$, all vertices except the last two can be colored sequentially using colors $\{1,2,3\}$ and the last two vertices can be colored using 4 and 5.\par

When $n \equiv 2 \text{ (mod 6)}$,  the colors $\{1,2,3,4\}$ can be sequentially used 5 times to color the first 20 vertices (leaving $n=14$ as a special case) and the remaining vertices, if any, can be colored sequentially using the colors $\{1,2,3\}$. When $n=14$, since a color can be used at most 3 times $(
\lfloor\frac{14}{3}\rfloor - 1 = 3)$ to color $C_{14}$, we need at least 5 colors to color $C_{14}$. But, 14 being even, and by Observation \ref{H odd}, the number of colors used to color $C_{14}$ must be even. Therefore, we need at least 6 colors to color $C_{14}$. Now, the colors $\{1,2,3,4\}$ sequentially used 3 times followed by 5 and 6, gives the required coloring of $C_{14}$. \par

When $n \equiv 3 \text{ (mod 6)}$, the colors $\{1,2,3\}$ can be used sequentially $\frac{n}{3}$ times to color the vertices of $C_n$. When When $n \equiv 4 \text{ (mod 6)}$, then again, the colors $\{1,2,3\}$ can be used sequentially $\lfloor\frac{n}{3}\rfloor$ times to color the $n-1$ vertices of $C_n$ and the last vertex can be colored using 4.\par

In all these cases, $K_1$, being a universal vertex, needs a color different from the colors used to color $C_n$. Hence, the theorem.
\end{proof}

\end{document}
